\section{Positivity and the closed quadratic form}\label{sec:positivity}

\subsection{Interval variables and the Bañados functional}\label{subsec:banados-functional}

Map $A=(-a,a)$ to the positive half-line and then to the real line by

\begin{align}
U(\phi)
=\frac{\sin\frac{\phi+a}{2}}
{\sin\frac{a-\phi}{2}},
\qquad
s=\log U.
\end{align}

The modular weight is

\begin{align}
w_A(\phi)=\frac{U}{U'}
=\frac{\cos\phi-\cos a}{\sin a}>0.
\end{align}

For $f_0=P_{\rm PSL}f$, define

\begin{align}
F_f(s)=\frac{f_0(\phi(s))}{w_A(\phi(s))}.
\end{align}

On the compact-anchor core, a finite chiral reparametrization can be written as $\rho(e^s)=e^{\varphi(s)}$. The vacuum Brown--Henneaux stress and the interval modular charge are

\begin{align}
\langle T_{UU}\rangle_\rho
=-\frac{c}{24\pi}\{\rho,U\},
\qquad
K_A=2\pi\int_0^\infty U T_{UU}\,\mathrm dU,
\end{align}

with $c=3/(2G)$. The Schwarzian chain rule gives

\begin{align}
U^2\{\rho,U\}
&=\{\varphi,s\}-\frac12\varphi'^2+\frac12,\\
\{\varphi,s\}
&=\partial_s\left(\frac{\varphi''}{\varphi'}\right)
-\frac12\left(\frac{\varphi''}{\varphi'}\right)^2.
\end{align}

The total derivative vanishes on the core. The charge-minus-geodesic functional is therefore

\begin{align}
\mathcal R_A[\rho]
=\frac{c}{24}\int_{\mathbb R}
\left[
\varphi'^2+\left(\frac{\varphi''}{\varphi'}\right)^2-1
\right]\mathrm ds . \label{eq:schwarzian-functional}
\end{align}

This is a classical Bañados/coadjoint-orbit computation. Its coefficient agrees with the chiral relative-entropy Schwarzian functional of Ref. \cite{Hollands:2019hje}, but no quantum relative-entropy theorem is assumed here.

For

\begin{align}
\varphi_\lambda(s)=s+\lambda F_f(s)+O(\lambda^2),
\end{align}

the Hessian of (\ref{eq:schwarzian-functional}) is

\begin{align}
\boxed{
E_{{\rm can,p}}[h[f]]
=\frac1{8G}\int_{\mathbb R}
\left[(F_f')^2+(F_f'')^2\right]\mathrm ds .
} \label{eq:positive-energy}
\end{align}

In interval variables the same form is

\begin{align}
8G E_{{\rm can,p}}[h[f]]
=\int_{-a}^{a}w_A
\left[
\left(\partial_\phi\frac{f_0}{w_A}\right)^2
+\left(
\partial_\phi\left[
w_A\partial_\phi\frac{f_0}{w_A}
\right]
\right)^2
\right]\mathrm d\phi . \label{eq:interval-energy}
\end{align}

\subsection{Kernel and positivity}\label{subsec:kernel-positivity}

Both terms in (\ref{eq:positive-energy}) are squares. Equality implies $F_f$ is constant, hence $f_0=Cw_A$. Restoring the two global directions removed by $P_{\rm PSL}$ gives

\begin{align}
\ker E_{{\rm can,p}}
=\operatorname{span}\{1,\cos\phi,\sin\phi\}
=\mathfrak{sl}(2,\mathbb R). \label{eq:energy-kernel}
\end{align}

These directions create no metric perturbation. Therefore the completed canonical energy is strictly positive on

\begin{align}
H^\sigma(S^1)/\mathfrak{sl}(2,\mathbb R).
\end{align}

There is no additional degeneracy in the charged $m\geq2$ sector. Positivity belongs to the complete form, not to the local $\Upsilon$ matrix by itself.

\subsection{Form domain and optimal ordinary regularity}\label{subsec:form-domain}

The intrinsic form domain is

\begin{align}
\mathscr D_A
=\left\{
[f]\in H^{1/2+}(S^1)/\mathfrak{sl}(2,\mathbb R):
F_f',F_f''\in L^2(\mathbb R,\mathrm ds)
\right\}. \label{eq:form-domain}
\end{align}

The compact-anchor core is dense in the norm induced by (\ref{eq:positive-energy}). Hence (\ref{eq:positive-energy}) is the unique positive closure of the finite-action/CPS form proved in Section \ref{sec:brown-henneaux-theorem}.

On the ordinary circle Sobolev scale,

\begin{align}
H^2(S^1)\cap\ker(q_+,q_-)
\hookrightarrow\mathscr D_A
\end{align}

continuously. Near an endpoint, use $y=w_A$ and write $f_0=yg$. Up to bounded smooth coefficients, the local energy is

\begin{align}
\int_0^{y_0}y
\left[g_y^2+(g_y+yg_{yy})^2\right]\mathrm dy .
\end{align}

The identity

\begin{align}
\int_0^{y_0}|2h+yh_y|^2\mathrm dy
=2\int_0^{y_0}|h|^2\mathrm dy
+\int_0^{y_0}y^2|h_y|^2\mathrm dy
+2y_0|h(y_0)|^2 \label{eq:hardy-identity}
\end{align}

with $h=g_y$ gives the required endpoint Hardy estimate.

The index $2$ is optimal. For an interior packet

\begin{align}
f_N(\phi)=N^{-r}\chi(\phi)e^{iN\phi}
\end{align}

with $\chi$ supported away from the anchors,

\begin{align}
\|f_N\|_{H^r}=O(1),
\qquad
E_{{\rm can,p}}[h[f_N]]\asymp N^{4-2r}.
\end{align}

Thus the form has no continuous extension to any ordinary $H^r$ with $r<2$.

It is not coercive in the $H^2$ norm. Let $\psi(x)=x^3(1-x)^3$ and place

\begin{align}
f_\delta(\phi)
=\delta^{3/2}\psi\left(\frac{\phi+a}{\delta}\right)
\end{align}

in an anchor layer. Then $\|f_\delta''\|_{L^2}$ approaches a nonzero constant while

\begin{align}
E_{{\rm can,p}}[h[f_\delta]]=O(\delta).
\end{align}

The completed energy is therefore positive but has no spectral gap in the ordinary $H^2$ topology.

At $H^2$, only the combined closed form is asserted. The separated pieces $E_{{\rm can,p}}[h]$ and $\int\Upsilon_{\rm p}$ can fail to exist individually because they contain sharper endpoint traces; their sum is the well-defined positive form.
